\sectionkicker{Source-backed technical notes}
\subsection*{生成から行動へ――RAG、Tool Use、Computer Use: Technical Notes}
本欄は記事本文で圧縮した一次資料上の情報を、比較・再検証しやすい形で整理したものである。時系列、確認済みの主張、留意点、一次資料URLを掲載する。
\smallskip
\noindent{\small\textit{Technical Notesでは、一次資料から確認した公開・提供・研究上の事実を記載する。数値・能力評価や提供元・著者による比較表現は、独立再現済みの結果ではなく帰属claimとして扱う。}}\par
\medskip
\subsection*{Theme at a glance}
\begin{center}
\footnotesize
\begin{tabularx}{\linewidth}{@{}p{0.20\linewidth}p{0.10\linewidth}p{0.14\linewidth}X@{}}
\toprule
資料 & 位置づけ & 種別 & 時系列 \\
\midrule
Command R / Command R+ & 主要資料 & モデル更新 & 2024-03-11, 2024-03-14, 2024-04-04           \\
SWE-agent & 主要資料 & Agent & 2024-05-06           \\
OSWorld & 主要資料 & 評価ベンチマーク & 2024-04-11           \\
OpenAI Model Spec & 補足資料 & 公式情報 & 2024-05-08           \\
Quiet-STaR & 補足資料 & 論文 & 2024-03-14           \\
\bottomrule
\end{tabularx}
\end{center}
\normalsize

\Needspace{0.38\textheight}
\begin{technicalnote}{Command R / Command R+}{主要資料}
\begingroup
% reader-facing Technical Notes break policy
\widowpenalty=10000
\clubpenalty=10000
\displaywidowpenalty=10000
\begin{tabularx}{\linewidth}{@{}>{\bfseries}p{0.22\linewidth}X@{}}
組織 & Cohere \\
種別 & モデル更新 \\
時系列 & 2024-03-11, 2024-03-14, 2024-04-04 \\
\end{tabularx}
\smallskip
\begin{minipage}{\linewidth}
% reader-facing Technical Notes event-only fact/source tail group
{\bfseries 一次資料から整理したtechnical points}
\begin{itemize}[leftmargin=1.5em,itemsep=0.35em]
\item \textbf{一次情報で確認できる事実}: Cohereの一次資料により、Command R / Command R+について2024-03-11にモデル公開、2024-03-14にTool Use発表、2024-04-04にDeployment提供開始を確認した。 Command Rは2024-03-11にRetrieval-Augmented Generationをproduction scaleで扱うモデルとして公開され、2024-03-14のTool Use発表ではCommand Rがuser-defined toolsと連携する利用形態が示された。Command R+は2024-04-04にenterprise-grade workloads向けのRAG-optimized modelとして公開された。
\end{itemize}
\begin{minipage}{\linewidth}
% reader-facing Technical Notes generic boundary/source tail group
\begin{samepage}
{\bfseries 一次資料}
\begingroup\sloppy
\Urlmuskip=0mu plus 2mu\relax
\begin{itemize}[leftmargin=1.5em,itemsep=0.25em]
\item {\scriptsize\url{https://cohere.com/blog/command-r}}
\item {\scriptsize\url{https://cohere.com/blog/command-r-plus-microsoft-azure}}
\item {\scriptsize\url{https://cohere.com/blog/tool-use-with-command-r}}
\end{itemize}
\endgroup
\end{samepage}
\end{minipage}
\end{minipage}
\endgroup
\end{technicalnote}
\medskip

\Needspace{0.38\textheight}
\begin{technicalnote}{SWE-agent}{主要資料}
\begingroup
% reader-facing Technical Notes break policy
\widowpenalty=10000
\clubpenalty=10000
\displaywidowpenalty=10000
\begin{tabularx}{\linewidth}{@{}>{\bfseries}p{0.22\linewidth}X@{}}
組織 & - \\
種別 & Agent \\
時系列 & 2024-05-06 \\
\end{tabularx}
\smallskip
\begin{minipage}{\linewidth}
% reader-facing Technical Notes event-only fact/source tail group
{\bfseries 一次資料から整理したtechnical points}
\begin{itemize}[leftmargin=1.5em,itemsep=0.35em]
\item \textbf{一次情報で確認できる事実}: 提供元の一次資料により、SWE-agentについて2024-05-06にPaper First Submissionを確認した。 SWE-agentはlanguage modelがsoftware repositoryを操作してGitHub issueを解く際のAgent-Computer Interface (ACI)を設計対象とし、repository navigation・file inspection・editing等の専用command/interfaceを与えてaction spaceとfeedbackを整える。SWE-bench上でissue resolutionを評価しているが、解決率はmodel・task subset・evaluation条件に依存する著者報告として扱う。
\end{itemize}
\begin{minipage}{\linewidth}
% reader-facing Technical Notes generic boundary/source tail group
\begin{samepage}
{\bfseries 一次資料}
\begingroup\sloppy
\Urlmuskip=0mu plus 2mu\relax
\begin{itemize}[leftmargin=1.5em,itemsep=0.25em]
\item {\scriptsize\url{http://arxiv.org/abs/2405.15793v3}}
\end{itemize}
\endgroup
\end{samepage}
\end{minipage}
\end{minipage}
\endgroup
\end{technicalnote}
\medskip

\Needspace{0.38\textheight}
\begin{technicalnote}{OSWorld}{主要資料}
\begingroup
% reader-facing Technical Notes break policy
\widowpenalty=10000
\clubpenalty=10000
\displaywidowpenalty=10000
\begin{tabularx}{\linewidth}{@{}>{\bfseries}p{0.22\linewidth}X@{}}
組織 & - \\
種別 & 評価ベンチマーク \\
時系列 & 2024-04-11 \\
\end{tabularx}
\smallskip
\begin{minipage}{\linewidth}
% reader-facing Technical Notes event-only fact/source tail group
{\bfseries 一次資料から整理したtechnical points}
\begin{itemize}[leftmargin=1.5em,itemsep=0.35em]
\item \textbf{一次情報で確認できる事実}: 提供元の一次資料により、OSWorldについて2024-04-11にPaper First Submissionを確認した。 OSWorldはUbuntu・Windows・macOS上の実際のcomputer environmentを対象とし、web/desktop applicationやOS操作をまたぐ369件のcomputer taskを、初期状態・自然言語instruction・実行後のstate-based evaluationで評価するbenchmarkとして構築された。agentはscreen/imageやaccessibility由来の観測からmouse/keyboard等のcomputer actionを実行し、成功率はauthors' evaluationとして扱う。
\end{itemize}
\begin{minipage}{\linewidth}
% reader-facing Technical Notes generic boundary/source tail group
\begin{samepage}
{\bfseries 一次資料}
\begingroup\sloppy
\Urlmuskip=0mu plus 2mu\relax
\begin{itemize}[leftmargin=1.5em,itemsep=0.25em]
\item {\scriptsize\url{http://arxiv.org/abs/2404.07972v2}}
\end{itemize}
\endgroup
\end{samepage}
\end{minipage}
\end{minipage}
\endgroup
\end{technicalnote}
\medskip

\Needspace{0.38\textheight}
\begin{technicalnote}{OpenAI Model Spec}{補足資料}
\begingroup
% reader-facing Technical Notes break policy
\widowpenalty=10000
\clubpenalty=10000
\displaywidowpenalty=10000
\begin{tabularx}{\linewidth}{@{}>{\bfseries}p{0.22\linewidth}X@{}}
組織 & OpenAI \\
種別 & 公式情報 \\
時系列 & 2024-05-08 \\
\end{tabularx}
\smallskip
\begin{minipage}{\linewidth}
% reader-facing Technical Notes event-only fact/source tail group
{\bfseries 一次資料から整理したtechnical points}
\begin{itemize}[leftmargin=1.5em,itemsep=0.35em]
\item \textbf{一次情報で確認できる事実}: OpenAIの一次資料により、OpenAI Model Specについて2024-05-08にSpec公開を確認した。 2024年5月8日のModel Spec初版は、望ましいmodel behaviorをObjectives・Rules・Default behaviorsの3層で整理し、Rulesで安全性・合法性に関わる境界を定めつつ、Defaultsはその境界内でdeveloper/userが上書き可能なstarting pointとして扱う構造を示した。OpenAIは当時、この文書をRLHFに関わるresearchers/data labelers向けのguidelineとして用いる意図を示す一方、現行形そのものはまだ学習に直接使っていないと明記しており、仕様公開と実モデル挙動を同一視しない。
\end{itemize}
\begin{minipage}{\linewidth}
% reader-facing Technical Notes generic boundary/source tail group
\begin{samepage}
{\bfseries 一次資料}
\begingroup\sloppy
\Urlmuskip=0mu plus 2mu\relax
\begin{itemize}[leftmargin=1.5em,itemsep=0.25em]
\item {\scriptsize\url{https://openai.com/index/introducing-the-model-spec}}
\end{itemize}
\endgroup
\end{samepage}
\end{minipage}
\end{minipage}
\endgroup
\end{technicalnote}
\medskip

\Needspace{0.38\textheight}
\begin{technicalnote}{Quiet-STaR}{補足資料}
\begingroup
% reader-facing Technical Notes break policy
\widowpenalty=10000
\clubpenalty=10000
\displaywidowpenalty=10000
\begin{tabularx}{\linewidth}{@{}>{\bfseries}p{0.22\linewidth}X@{}}
組織 & - \\
種別 & 論文 \\
時系列 & 2024-03-14 \\
\end{tabularx}
\smallskip
\begin{minipage}{\linewidth}
% reader-facing Technical Notes event-only fact/source tail group
{\bfseries 一次資料から整理したtechnical points}
\begin{itemize}[leftmargin=1.5em,itemsep=0.35em]
\item \textbf{一次情報で確認できる事実}: 提供元の一次資料により、Quiet-STaRについて2024-03-14にPaper First Submissionを確認した。 Quiet-STaRは、将来のテキストを説明するrationaleを各token位置で生成する学習法として、tokenwise parallel sampling、thoughtの開始・終了を示すlearnable token、extended teacher forcingを組み合わせる。著者らはcontinued pretraining後のGSM8K・CommonsenseQA等の改善を報告しているが、数値性能は著者報告として扱う。
\end{itemize}
\begin{minipage}{\linewidth}
% reader-facing Technical Notes generic boundary/source tail group
\begin{samepage}
{\bfseries 一次資料}
\begingroup\sloppy
\Urlmuskip=0mu plus 2mu\relax
\begin{itemize}[leftmargin=1.5em,itemsep=0.25em]
\item {\scriptsize\url{http://arxiv.org/abs/2403.09629v2}}
\end{itemize}
\endgroup
\end{samepage}
\end{minipage}
\end{minipage}
\endgroup
\end{technicalnote}
\medskip
